\documentclass[12pt]{article}
\usepackage[utf8]{inputenc}
\usepackage[T1]{fontenc}
\usepackage{hyperref}
\title{Tedaga Language Resources for K-12 Curriculum Development}
\author{Compiled by Claude Sonnet 5 (claude-sonnet-5) for Andr\'as Kornai}
\date{}
\begin{document}
\maketitle

\section{Introduction}
Tedaga (Teda, ISO 639-3 tuq) is a Saharan (Nilo-Saharan) language of the northern Tubu (Toubou) people in northern Chad, southern Libya and eastern Niger. It is closely related to Daza, and the two together make up Tebu. A Latin-based alphabet is in use. Web search turned up little material for schools. Most of what exists is older descriptive work, plus recent community dictionary projects. No Tedaga curriculum or school programme was found.

\section{Primary Educational Resources}
No primers, textbooks or school programmes for Tedaga turned up in search. The community dictionary apps below could serve as vocabulary support for classroom materials.

\section{Linguistic and Reference Materials}
\begin{itemize}
\item Teda language, Wikipedia overview. Its references list the main older works: Le C\oe{}ur, \emph{Dictionnaire ethnographique t\'eda} (1950); Le C\oe{}ur and Le C\oe{}ur, \emph{Grammaire et textes teda-daza} (1956); Jourdan (1935); Lukas (1953). \url{https://en.wikipedia.org/wiki/Teda_language}
\item Webonary online Tudaga dictionary from the Mosko Hanadii-\~i project in Bardai, Chad. It is multilingual (Tudaga, French, Arabic, English). \url{https://www.webonary.org/tudaga/overview/introduction/?lang=en}
\item Omniglot chart for Teda. \url{https://www.omniglot.com/charts/teda.pdf}
\end{itemize}
The Haggar and Walters \emph{Lexique tubu (dazaga)-fran\c{c}ais} (SIL Niger, 2005) is reported to cover the closely related Daza variety. I found no link to it.

\section{Religious and Community Resources}
None found in search.

\section{Digital and Audio Resources}
\begin{itemize}
\item Mobile dictionary app ``M\^ede t\^ugi tudagaa'' (Tudaga to Arabic, English and French, and back). Its own description calls it preliminary. \url{https://apps.apple.com/us/app/id1536571686}
\item Tedaga W\"orterbuch, a German-linked dictionary project with roughly 3,000 entries. \url{https://www.fundraiso.ch/en/organisations/tedaga-worterbuch}
\end{itemize}

\section{Curriculum Development Context}
Tedaga is not known to be taught in schools. A curriculum would have to start from scratch: standardise the orthography with community input, build on the Webonary dictionary, and write primers. Speaker counts differ across sources (about 29,000 in the survey data, higher elsewhere). Daza materials may be partly reusable, but this needs confirming with speakers.
\end{document}